\section{Cómputo y restricciones organizacionales: el cuello de botella real de la ciencia abierta}

\subsection{El cómputo sigue siendo la restricción de primer orden}
El ponente señala que, incluso con una definición clara de apertura y una comunidad con voluntad firme, la ciencia abierta seguirá estando limitada por el cómputo y por la ingeniería de sistemas.
Esto significa que la competencia futura no solo estará en la capacidad de los modelos, sino también en quién pueda construir infraestructura de entrenamiento y evaluación reutilizable.

\begin{figure}[H]
\centering
\includegraphics[width=0.93\textwidth]{frame-07-adoption.jpg}
\caption{Adopción del ecosistema y contribución de la comunidad: los efectos de red de los marcos abiertos}
\end{figure}
\noindent\footnotesize Evidencia en pantalla: subtítulos aproximadamente en 00:26:01 y 00:31:01, donde el ponente aborda la escala del ecosistema, las prácticas de los modelos abiertos y el valor de la migración de la investigación.\normalsize

\begin{knowledgebox}{Tres tipos de recursos escasos desde la perspectiva de la infraestructura}
\begin{itemize}
    \item Cómputo de entrenamiento: determina el espacio de modelos y datos que se puede explorar;
    \item Talento de ingeniería: determina la mantenibilidad del sistema y el rendimiento de los experimentos;
    \item Activos de evaluación: determinan si el progreso se puede verificar y comparar.
\end{itemize}
\end{knowledgebox}

\subsection{Una vía viable para la colaboración descentralizada}
\begin{importantbox}{La verdadera palanca del ecosistema abierto: estándares y cadenas de herramientas compartidos}
El cómputo de un solo punto difícilmente puede replicar el entrenamiento frontier, pero la comunidad puede formar un «multiplicador de colaboración» mediante evaluaciones unificadas, scripts de entrenamiento portables y procesos reutilizables de gobernanza de datos.
\end{importantbox}

\begin{warningbox}{Sin estandarización, la colaboración degenera en trabajo repetido y fragmentado}
Si cada equipo usa formatos de registro, protocolos de evaluación y definiciones de tareas incompatibles entre sí, el resultado es que «todos están haciendo experimentos, pero nadie puede compararlos de manera efectiva».
\end{warningbox}

\subsection{Resumen de la sección}
La ciencia abierta no equivale a una barrera de entrada baja. Requiere nueva infraestructura común para reducir la fricción de la colaboración; de lo contrario, la «apertura» se quedará solo en el nivel del lema.
